\section{Introduction}
\label{sec:introduction}

Multi-horizon time series forecasting produces predictions for several future time steps from the observations available at a forecast origin. As new observations arrive, the errors of earlier forecasts become available. These residuals can contain information that the original predictor has not captured. Residual correction methods have used this information to improve subsequent predictions in traffic forecasting and other time series tasks~\cite{rescal2022,tefl2026}.

We consider a setting in which the parameters of a trained base forecaster remain fixed during evaluation. An auxiliary correction model uses completed forecast residuals to adjust later forecasts. This separates adaptation of the correction from retraining of the base model. The objective is to improve forecasting accuracy while limiting the numerical state retained by the correction model between predictions. Both quantities matter when comparing correction methods, because a reduction in error may require a larger regression model or additional stored observations.

Several existing approaches provide a basis for this objective. ResCAL estimates future errors from previous errors and graph signals~\cite{rescal2022}. Temporal Error Feedback Learning (TEFL) incorporates historical residuals through a learned low-rank adapter and jointly trains the adapter with the base forecaster~\cite{tefl2026}. ELF combines an online forecaster with a fixed forecasting model through adaptive weighting~\cite{elf2025}. These approaches establish that feedback can be incorporated through residual estimation or an auxiliary online predictor. Our focus is the representation supplied to a correction model when the retained state is compressed.

A fixed low-dimensional projection provides a simple way to summarize a residual block. However, the projection generally does not preserve each individual residual value. In particular, two blocks can have the same retained coefficients but different residuals at the final time step. A regression model that receives only those coefficients cannot distinguish the two blocks. The final residual is also the most recent error in a fully observed block. This motivates testing whether retaining it explicitly improves correction beyond what can be obtained from the compressed coefficients alone.

We propose endpoint-preserving compressed residual correction. Here, the endpoint denotes the residual at the final time step of the preceding fully observed forecast block. The method retains fixed discrete cosine transform (DCT) coefficients of that block together with its uncompressed endpoint residual. It combines these quantities with DCT coefficients of the current base forecast in a separate online ridge regression for each retained component and channel. The estimated correction is reconstructed in the same DCT subspace and blended with the base forecast. Updates use only forecast blocks whose target values have become available. The proposed contribution is this state and input construction. Ridge regression, the DCT, and forecast blending provide established components for its implementation.

The evaluation distinguishes the effect of endpoint selection from the performance of the complete correction method. Comparisons with the first residual, a middle residual, and the block mean hold the state size and regression procedure fixed. Separate comparisons use learned residual adapters and implementations of ELF adapted to the same forecasting protocol. On the evaluated real-data conditions, endpoint retention improves mean squared error (MSE) and mean absolute error (MAE) over the alternative summaries. The complete method also improves accuracy over the evaluated learned adapters. These results do not imply simultaneous superiority in accuracy and storage over every comparator. A larger ELF configuration achieves lower average error, whereas smaller learned adapters retain fewer bytes. We therefore report accuracy and retained numerical storage separately, including conditions and intervals in which correction is less effective.

The contributions are as follows:
\begin{enumerate}
\item We introduce a compressed residual representation that explicitly retains the endpoint and uses the current base forecast as additional input to online correction. This construction preserves a residual value that is generally not recoverable from the retained DCT coefficients.
\item We evaluate the endpoint choice using alternative summaries with the same state size. These comparisons isolate its contribution from changes in regression capacity and the update procedure.
\item We characterize the accuracy--storage tradeoff against learned residual adapters and online Fourier regression. The evaluation accounts for numerical storage under comparable representations and identifies both average improvements and failure cases.
\end{enumerate}

Section~\ref{sec:related_work} discusses related methods. Section~\ref{sec:method} defines the proposed correction, and Section~\ref{sec:experimental_setup} specifies the evaluation protocol. Section~\ref{sec:results} presents the results. Sections~\ref{sec:discussion} and~\ref{sec:conclusion} discuss their scope and conclude the paper.
